\documentclass{article}

\usepackage[utf8]{inputenc}
\usepackage[T1]{fontenc}
\usepackage{helvet}
\usepackage{geometry}
\usepackage{xcolor}
\usepackage{eso-pic}
\usepackage{fancyhdr}
\usepackage{parskip}
\usepackage{microtype}
\usepackage[english, ukrainian]{babel}
\usepackage{url}
\usepackage{amsmath}
\usepackage{amssymb}
\usepackage{amsthm}
\usepackage{longtable}
\usepackage{booktabs}
\usepackage{hyperref}
\usepackage{titlesec}
\usepackage{tocloft}

\renewcommand{\cftsecfont}{\small\bfseries}
\renewcommand{\cftsubsecfont}{\footnotesize}
\renewcommand{\cftsubsubsecfont}{\scriptsize}
\renewcommand{\cftsecpagefont}{\small\bfseries}
\renewcommand{\cftsubsecpagefont}{\footnotesize}
\renewcommand{\cftsubsubsecpagefont}{\scriptsize}

\definecolor{nato}{HTML}{293757}
\definecolor{footergray}{gray}{0.65}

\geometry{a4paper,left=3cm,right=3cm,top=3cm,bottom=3cm}
\renewcommand{\familydefault}{\sfdefault}
\setcounter{secnumdepth}{3}

\titleformat{\section}
  {\color{nato}\Huge\bfseries\raggedright}{\thesection.}{1em}{}
\titlespacing*{\section}{0pt}{30pt}{12pt}

\titleformat{\subsection}
  {\color{nato}\LARGE\bfseries\raggedright}{\thesubsection.}{1em}{}
\titlespacing*{\subsection}{0pt}{20pt}{8pt}

\titleformat{\subsubsection}
  {\color{nato}\Large\bfseries\raggedright}{\thesubsubsection.}{1em}{}
\titlespacing*{\subsubsection}{0pt}{10pt}{6pt}

\fancyhf{}
\fancyhead[R]{\small\color{footergray} 2 червня 2026}
\fancyfoot[L]{\small\color{footergray} Комплексна система захисту інформації. Галузевий профіль для месенджера (20).}
\fancyfoot[R]{\small\color{footergray}\thepage}

\newcommand\HUGE[1]{\fontsize{40}{40}\selectfont #1}

\renewcommand{\headrulewidth}{0pt}
\renewcommand{\footrulewidth}{0.4pt}
\renewcommand{\footrule}{\color{footergray}\hrule width \textwidth height 0.4pt}

\begin{document}

\pagestyle{fancy}

\begin{titlepage}
  \AddToShipoutPictureBG*{\AtPageLowerLeft{\color{nato}\rule{\paperwidth}{\paperheight}}}
  \color{white}\thispagestyle{empty}\vspace*{4cm}\centering
  {\Huge  \textbf{Комплексна система захисту інформації} \par} \vspace{2cm}
  {\HUGE  \textbf{Галузевий профіль для месенджера (20)} \par} \vspace{2.5cm}
  \vfill
  {\large 2026 \copyright\ Критпографічні Телесистеми \par}
\end{titlepage}

\newpage

\begin{center}
  {\Huge\color{nato}\bfseries Зміст\par}
\end{center}
\vspace{40pt}
\tableofcontents

\begin{abstract}
Галузевий профіль — розробляється для забезпечення конфіденційності, цілісності повідомлень та управління сесіями в месенджері.
\end{abstract}

\section{AC}
Клас заходів захисту AC --- УПРАВЛІННЯ ДОСТУПОМ

Цей клас зосереджується на обмеженні доступу до інформаційних систем, ресурсів та функцій лише для авторизованих користувачів, програм і пристроїв.

\textbf{Перелік заходів захисту:} Управління паралельною сесією (AC-10); Id-spe-ac-19-1; Id-spe-ac-19-5.

\subsection{УПРАВЛІННЯ ПАРАЛЕЛЬНОЮ СЕСІЄЮ (AC-10)}
Обмежити кількість одночасних сеансів для кожного [Призначення: визначеного організацією облікового запису та/або типу облікового запису] до [Призначення: визначеної організацією кількості].

\vspace{10pt}
{\footnotesize\bfseries
No: 1\\
Name: ac\_10\_01\\
Type: integer\\
Default: 30
}

{\footnotesize Кількість одночасних сеансів для кожного облікового запису та/або типів облікових записів обмежена кількість}


{\footnotesize\bfseries
No: 2\\
Name: ac\_10\_odp\_02\\
Type: integer\\
Default: 30
}

{\footnotesize Визначено кількість одночасних сеансів, дозволених для кожного облікового запису та/або типу облікового запису}




\subsubsection{id-spe-ac-19-1}


\vspace{10pt}
{\footnotesize\bfseries
No: 1\\
Name: ac\_19\_1\_01\\
Type: string\\
Default: nil
}

{\footnotesize КОНТРОЛЬ ДОСТУПУ ДЛЯ МОБІЛЬНИХ ПРИСТРОЇВ - ВИКОРИСТАННЯ ПИСЬМОВИХ ТА ПОРТАТИВНИЙ ПРИСТРОЇВ ДЛЯ ЗБЕРІГАННЯ ДАНИХ [Вилучено: Включено в MP-07]}




\subsubsection{id-spe-ac-19-5}


\vspace{10pt}
\textit{Немає параметрів для цього контролю.}
\vspace{10pt}


\section{AU}
Клас заходів захисту AU --- АУДИТ ТА ПІДЗВІТНІСТЬ

Цей клас забезпечує можливість відстеження дій у системі, генерування, захист та аналіз записів аудиту для виявлення порушень політики безпеки.

\textbf{Перелік заходів захисту:} Неспростовність (AU-10); Id-spe-au-10-5.

\subsection{НЕСПРОСТОВНІСТЬ (AU-10)}
Надавайте неспростовні докази того, що особа (або процес, який діє від імені особи) виконала [Призначення: дії, визначені організацією, на які поширюється принцип неспростовності].

\vspace{10pt}
{\footnotesize\bfseries
No: 1\\
Name: au\_10\_01\\
Type: list\\
Default: [\textquotedbl{}admin\textquotedbl{}, \textquotedbl{}security\_officer\textquotedbl{}]
}

{\footnotesize Надаються неспростовні докази того, що особа (або процес, що діє від імені особи) виконала дії}


{\footnotesize\bfseries
No: 2\\
Name: au\_10\_odp\\
Type: list\\
Default: [\textquotedbl{}login\textquotedbl{}, \textquotedbl{}logout\textquotedbl{}, \textquotedbl{}failed\_attempt\textquotedbl{}]
}

{\footnotesize Визначено дії, на які поширюється принцип неспростовності}




\subsubsection{id-spe-au-10-5}


\vspace{10pt}
{\footnotesize\bfseries
No: 1\\
Name: au\_10\_5\_01\\
Type: string\\
Default: nil
}

{\footnotesize НЕСПРОСТОВНІСТЬ - ЦИФРОВІ ПІДПИСИ [Вилучено: Включено до SI-07]}




\section{IA}
Клас заходів захисту IA --- ІДЕНТИФІКАЦІЯ ТА АВТЕНТИФІКАЦІЯ

Цей клас відповідає за однозначне розпізнавання користувачів або пристроїв та підтвердження їхньої справжності перед наданням доступу до системи.

\textbf{Перелік заходів захисту:} Id-spe-ia-5-12; Id-spe-ia-7.

\subsubsection{id-spe-ia-5-12}


\vspace{10pt}
{\footnotesize\bfseries
No: 1\\
Name: ia\_5\_12\_odp\\
Type: string\\
Default: \textquotedbl{}автоматизований засіб моніторингу\textquotedbl{}
}

{\footnotesize Визначено вимоги до якості біометрії; для біометричної автентифікації використовувати механізми, які задовольняють вимоги}




\subsection{id-spe-ia-7}


\vspace{10pt}
{\footnotesize\bfseries
No: 1\\
Name: ia\_7\_01\\
Type: list\\
Default: [\textquotedbl{}default\_deny\_rule\textquotedbl{}, \textquotedbl{}abac\_rule\_1\textquotedbl{}]
}

{\footnotesize Впроваджено механізми автентифікації в криптографічний модуль, який відповідає вимогам чинних законів, виконавчих розпоряджень, директив, політик, правил, стандартів та рекомендацій для такої автентифікації}




\section{MP}
Клас заходів захисту MP --- ЗАХИСТ НОСІЇВ ІНФОРМАЦІЇ

Цей клас спрямований на безпечне зберігання, транспортування, використання та знищення як цифрових, так і паперових носіїв інформації.

\textbf{Перелік заходів захисту:} Id-spe-mp-6.

\subsection{id-spe-mp-6}
a. Очищувати [Призначення: визначені організацією системні носії] перед утилізацією, випуском за межі організаційного контролю, або перед повторним використанням [Призначення: методами та процедурами очищення, визначеними організацією].\\ 
b. Використовувати механізми очищення зі стійкістю та цілісністю, що відповідає категорії безпеки або рівню секретності інформації.

\vspace{10pt}
{\footnotesize\bfseries
No: 1\\
Name: mp\_6\_a\_01\\
Type: string\\
Default: nil
}

{\footnotesize Носії інформації системи перед утилізацією піддаються очищенню за допомогою методів та процедур очищення}


{\footnotesize\bfseries
No: 2\\
Name: mp\_6\_a\_02\\
Type: string\\
Default: nil
}

{\footnotesize Носії інформації системи очищаються за допомогою методів та процедур очищення перед випуском за межі контрольованої зони}


{\footnotesize\bfseries
No: 3\\
Name: mp\_6\_a\_03\\
Type: string\\
Default: nil
}

{\footnotesize Носії інформації системи очищуються за допомогою методів і процедур очищення перед повторним використанням}


{\footnotesize\bfseries
No: 4\\
Name: mp\_6\_b\\
Type: string\\
Default: \textquotedbl{}автоматизований засіб моніторингу\textquotedbl{}
}

{\footnotesize Застосовуються механізми очищення, надійність і цілісність яких відповідає категорії безпеки або рівню секретності інформації}


{\footnotesize\bfseries
No: 5\\
Name: mp\_6\_odp\_01\\
Type: string\\
Default: nil
}

{\footnotesize Визначено носії інформації системи, які підлягають очищенню перед утилізацією}


{\footnotesize\bfseries
No: 6\\
Name: mp\_6\_odp\_02\\
Type: string\\
Default: nil
}

{\footnotesize Визначено носії інформації системи, які підлягають очищенню перед випуском за межі контрольованої зони}


{\footnotesize\bfseries
No: 7\\
Name: mp\_6\_odp\_03\\
Type: string\\
Default: nil
}

{\footnotesize Визначені носії інформації системи, що підлягають очищенню перед повторним використанням}


{\footnotesize\bfseries
No: 8\\
Name: mp\_6\_odp\_04\\
Type: string\\
Default: nil
}

{\footnotesize Визначено методи та процедури очищення, які слід використовувати для очищення перед утилізацією}


{\footnotesize\bfseries
No: 9\\
Name: mp\_6\_odp\_05\\
Type: string\\
Default: nil
}

{\footnotesize Визначено методи та процедури очищення, які слід використовувати для очищення перед випуском за межі контрольованої зони}


{\footnotesize\bfseries
No: 10\\
Name: mp\_6\_odp\_06\\
Type: string\\
Default: nil
}

{\footnotesize Визначено методи та процедури очищення, які слід використовувати для очищення перед повторним використанням}




\section{SC}
Клас заходів захисту SC --- ЗАХИСТ ІНФОРМАЦІЙНОЇ СИСТЕМИ ТА

Цей клас охоплює механізми захисту інформації під час її передачі мережами зв'язку, криптографічний захист та ізоляцію критичних компонентів.

\textbf{Перелік заходів захисту:} Id-spe-sc-5; Id-spe-sc-8-1; Id-spe-sc-8-2; Встановлення ключами (SC-12); Id-spe-sc-17; Id-spe-sc-20; Id-spe-sc-21; Id-spe-sc-22; Автентифікація сесії (SC-23); Id-spe-sc-28-1; Id-spe-sc-49; Апаратний захист (SC-51).

\subsection{id-spe-sc-5}
a. [Призначення: захистити від; Обмежити] наслідки наступних типів подій відмови в обслуговуванні (DoS): [Призначення: визначені організацією типи подій відмови в обслуговуванні];\\ 
b. Застосувати наступні заходи захисту для досягнення мети відмови обслуговування [Призначення: заходи захисту визначені організацією, за типом події відмови в обслуговуванні].

\vspace{10pt}
{\footnotesize\bfseries
No: 1\\
Name: sc\_5\_odp\_01\\
Type: string\\
Default: nil
}

{\footnotesize Визначені типи подій відмов в обслуговуванні, від яких потрібно захищати або обмежувати; SC-05\_ODP[02] вибрано одне з наступних ЗНАЧЕНЬ ПАРАМЕТРІВ: \{захистити від; обмежити\}}




\subsubsection{id-spe-sc-8-1}


\vspace{10pt}
{\footnotesize\bfseries
No: 1\\
Name: sc\_8\_1\_odp\\
Type: string\\
Default: nil
}

{\footnotesize Вибрано одне або декілька з наступних ЗНАЧЕНЬ ПАРАМЕТРІВ: \{запобігти несанкціонованому розголошенню інформації; виявити зміни в інформації\}}




\subsubsection{id-spe-sc-8-2}


\vspace{10pt}
{\footnotesize\bfseries
No: 1\\
Name: sc\_8\_2\_01\\
Type: string\\
Default: nil
}

{\footnotesize КОНФІДЕНЦІЙНІСТЬ ТА ЦІЛІСНІСТЬ ПЕРЕДАЧІ - ПОПЕРЕДНЯ І ПОСТОБРОБКА МЕТА ОЦІНКИ: Визначити, чи:}


{\footnotesize\bfseries
No: 2\\
Name: sc\_8\_2\_odp\\
Type: string\\
Default: nil
}

{\footnotesize Вибрано одне або декілька з наступних ЗНАЧЕНЬ ПАРАМЕТРІВ: \{конфіденційність; цілісність\}}




\subsection{ВСТАНОВЛЕННЯ КЛЮЧАМИ (SC-12)}
Встановити та управляти криптографічними ключами для криптографічних засобів, які використовуються в системі відповідно до [Призначення: визначені організацією вимоги до генерації, поширення, зберігання, доступу та знищення ключів].

\vspace{10pt}
{\footnotesize\bfseries
No: 1\\
Name: sc\_12\_01\\
Type: string\\
Default: \textquotedbl{}AES-256-GCM\textquotedbl{}
}

{\footnotesize Встановлюються криптографічні ключі, коли в системі використовується криптографія відповідно до < SC-12\_ODP вимог >}


{\footnotesize\bfseries
No: 2\\
Name: sc\_12\_02\\
Type: string\\
Default: \textquotedbl{}AES-256-GCM\textquotedbl{}
}

{\footnotesize Здійснюється управління криптографічними ключами, коли в системі використовується криптографія, відповідно до вимог }


{\footnotesize\bfseries
No: 3\\
Name: sc\_12\_odp\\
Type: string\\
Default: nil
}

{\footnotesize Визначені вимоги до генерації, розповсюдження, зберігання, доступу та знищення ключів}




\subsection{id-spe-sc-17}
a. Випускати сертифікати відкритого ключа відповідно до [Призначення: визначеної організацією політики сертифікації];\\ 
b. Отримувати сертифікати відкритого ключа від затвердженого постачальника послуг.

\vspace{10pt}
\textit{Немає параметрів для цього контролю.}
\vspace{10pt}


\subsection{id-spe-sc-20}
a. Надати додаткові дані автентифікації та перевірки цілісності джерела даних разом з офіційними даними розпізнавання імен, які система повертає у відповідь на запити дозволу імен/адрес.\\ 
b. Надати засоби для вказання статусу безпеки дочірніх зон і (якщо дочірня зона підтримує служби безпечного дозволу) забезпечити перевірку ланцюга довіри між батьківськими та дочірніми доменами при роботі в складі розподіленого ієрархічного простору імен.

\vspace{10pt}
\textit{Немає параметрів для цього контролю.}
\vspace{10pt}


\subsection{id-spe-sc-21}
Зробити запит та виконати перевірку автентичності джерела даних і перевірку цілісності даних у відповідях на дозвіл імен/адрес, які система отримує від уповноважених джерел.

\vspace{10pt}
{\footnotesize\bfseries
No: 1\\
Name: sc\_21\_01\\
Type: string\\
Default: nil
}

{\footnotesize Реалізується запит перевірки автентичності джерела даних для відповідей на запит дозволу імен/адрес, які система отримує від авторитетних джерел}


{\footnotesize\bfseries
No: 2\\
Name: sc\_21\_02\\
Type: string\\
Default: nil
}

{\footnotesize Реалізується запит автентифікація походження даних на основі відповідей з дозволу імен/адрес, які система отримує від авторитетних джерел}


{\footnotesize\bfseries
No: 3\\
Name: sc\_21\_03\\
Type: string\\
Default: nil
}

{\footnotesize Реалізується запит перевірки цілісності даних для відповідей на запит дозволу імен/адрес, які система отримує від авторитетних джерел}


{\footnotesize\bfseries
No: 4\\
Name: sc\_21\_04\\
Type: string\\
Default: nil
}

{\footnotesize Виконується перевірка цілісності даних для відповідей на запит про дозвіл імен/адрес, які система отримує від авторитетних джерел}




\subsection{id-spe-sc-22}
Переконатися, що системи, які спільно надають послуги розпізнавання імен/адрес для організації, є відмовостійкими та забезпечують поділ внутрішніх і зовнішніх ролей.

\vspace{10pt}
{\footnotesize\bfseries
No: 1\\
Name: sc\_22\_01\\
Type: string\\
Default: nil
}

{\footnotesize Є системи, які спільно надають послуги з визначення імен/адрес для організації, відмовостійкими}


{\footnotesize\bfseries
No: 2\\
Name: sc\_22\_02\\
Type: string\\
Default: nil
}

{\footnotesize Реалізовано в системах, які спільно надають послуги з вирішення імен/адрес для організації, внутрішній розподіл ролей}


{\footnotesize\bfseries
No: 3\\
Name: sc\_22\_03\\
Type: string\\
Default: nil
}

{\footnotesize Реалізовано в системах, які спільно надають послуги з вирішення імен/адрес для організації, зовнішній розподіл ролей}




\subsection{АВТЕНТИФІКАЦІЯ СЕСІЇ (SC-23)}
Забезпечити автентифікацію сеансів зв'язку.

\vspace{10pt}
{\footnotesize\bfseries
No: 1\\
Name: sc\_23\_01\\
Type: string\\
Default: nil
}

{\footnotesize Захищено автентифікацію сеансів зв'язку}




\subsubsection{id-spe-sc-28-1}


\vspace{10pt}
{\footnotesize\bfseries
No: 1\\
Name: sc\_28\_1\_01\\
Type: string\\
Default: \textquotedbl{}AES-256-GCM\textquotedbl{}
}

{\footnotesize Реалізовані криптографічні механізми для запобігання несанкціонованому розкриттю інформації, що знаходиться в стані спокою на системних компонентах або носіях}


{\footnotesize\bfseries
No: 2\\
Name: sc\_28\_1\_02\\
Type: string\\
Default: \textquotedbl{}AES-256-GCM\textquotedbl{}
}

{\footnotesize Реалізовані криптографічні механізми для запобігання несанкціонованій модифікації інформації, що знаходиться в стані спокою на системних компонентах або носіях}




\subsection{id-spe-sc-49}
Впровадити механізми апаратного поділу та застосування політики між [Призначення: домени безпеки, визначені організацією].

\vspace{10pt}
{\footnotesize\bfseries
No: 1\\
Name: sc\_49\_01\\
Type: list\\
Default: [\textquotedbl{}default\_deny\_rule\textquotedbl{}, \textquotedbl{}abac\_rule\_1\textquotedbl{}]
}

{\footnotesize Впроваджено механізми апаратного розділення та застосування політик між доменами безпеки}


{\footnotesize\bfseries
No: 2\\
Name: sc\_49\_odp\\
Type: list\\
Default: [\textquotedbl{}default\_deny\_rule\textquotedbl{}, \textquotedbl{}abac\_rule\_1\textquotedbl{}]
}

{\footnotesize Визначені домени безпеки, які потребують апаратного розділення та механізмів забезпечення дотримання політики}




\subsection{АПАРАТНИЙ ЗАХИСТ (SC-51)}
a. Перевіряти правильність роботи [Призначення: визначені організацією функції безпеки та приватності].\\ 
b. Виконувати перевірку [Вибір (один або кілька): [Призначення: визначені організацією системні перехідні стани]; за командою користувача з відповідними повноваженнями; [Призначення: визначена організацією частота]].\\ 
c. Повідомляти [Призначення: визначені організацією персонал або посадові особи] про невдалі перевірки безпеки та приватності.\\ 
d. [Вибір (один або кілька): Вимкнути систему; Перезапустити систему; [Призначення: визначені організацією альтернативні дії]], коли виявляються аномалії.

\vspace{10pt}
{\footnotesize\bfseries
No: 1\\
Name: sc\_51\_odp\_01\\
Type: string\\
Default: nil
}

{\footnotesize Визначено компоненти системної прошивки, потребують апаратного захисту від запису; які}


{\footnotesize\bfseries
No: 2\\
Name: sc\_51\_odp\_02\\
Type: list\\
Default: [\textquotedbl{}admin\textquotedbl{}, \textquotedbl{}security\_officer\textquotedbl{}]
}

{\footnotesize Визначені уповноважені особи, які повинні виконувати процедури вимкнення та повторного ввімкнення апаратного захисту від запису; SC-51a. використовується апаратний захист від запису для компонентів мікропрограми системи; SC-51b.[01] впроваджено спеціальні процедури для уповноважених осіб для ручного вимкнення апаратного захисту від запису для модифікацій мікропрограми; SC-51b.[02] реалізовано спеціальні процедури для уповноважених осіб для повторного увімкнення захисту від запису перед поверненням до робочого режиму}




\end{document}
